\chapter{Latent Reduced-Order Model: MVAR and LSTMs in POD Coordinates}
\label{ch:latent_rom}

This chapter describes the two latent-space time-steppers used in the
restrict--evolve--lift pipeline of Chapter~\ref{ch:global_pod_pipeline}:
a multivariate autoregressive model~(MVAR) and a long short-term memory
network~(LSTM)~\cite{hochreiter1997lstm}.
Both models operate on the $r$-dimensional POD coefficient
vector~\cite{berkooz1993pod}
$y(t_k)\in\mathbb{R}^r$ and produce a one-step-ahead prediction that is
iterated in closed loop to generate long-horizon density forecasts.
The purpose of this chapter is not to search for the best black-box
predictor in the abstract.  Rather, it is to test whether shift alignment
and POD produce a latent representation whose time evolution is effectively
linear at forecast-relevant scales, or whether substantial nonlinear memory
remains.  If MVAR performs well, that suggests alignment has removed a major
source of apparent complexity; if LSTM materially outperforms MVAR, genuine
nonlinear latent structure persists; and if both fail, the issue may be
low-dimensionality itself rather than model class.

%% ================================================================
\section{MVAR(\texorpdfstring{$w$}{w}) model}
\label{sec:mvar_definition}
%% ================================================================

Let $x(t)$ denote the macroscopic density field on a fixed grid, sampled at
discrete times $t_k := t_0 + k\,\Delta t$ ($k = 0,1,2,\dots$) with effective sampling interval
$\Delta t = 0.04\times 3 = 0.12$\,s.
The restriction operator $\mathcal{R}$ (Chapter~\ref{ch:global_pod_pipeline}) maps
$x(t_k)$ to a latent coordinate $y(t_k)\in\mathbb{R}^r$ with
$r=19$.
In this chapter we focus on learning a surrogate time-stepper for $y$.

\subsection{Model class}
We model the latent coordinates as a discrete-time vector autoregression with
$w$ lags.
For each training trajectory (``case'') $c$ we write
\begin{equation}\label{eq:mvar_def}
  y^{(c)}(t_k) \;=\; A_0 \;+\; \sum_{j=1}^{w} A_j\,y^{(c)}(t_{k-j})
                 \;+\; \varepsilon(t_k),
\end{equation}
where $A_0\in\mathbb{R}^r$ is an intercept vector and
$A_j\in\mathbb{R}^{r\times r}$ are lag-$j$ regression matrices~\cite{alvarez2025eqfree}.
The residual $\varepsilon(t_k)\in\mathbb{R}^r$ is modelled as zero-mean
white noise; this assumption motivates least-squares training but is not
used for inference (we do point prediction).
MVAR serves here as a hypothesis test: after shift alignment and POD, is
the dominant latent evolution well approximated by a linear finite-memory
map?

\noindent\textbf{One-step predictor.}
The corresponding one-step-ahead prediction is
\begin{equation}\label{eq:mvar_predictor}
  \hat y^{(c)}(t_k) \;=\; A_0 \;+\;
  \sum_{j=1}^{w} A_j\,y^{(c)}(t_{k-j}).
\end{equation}

\subsection{Ridge-regularized training}\label{subsec:mvar_training}
The unknown parameters $(A_0,\{A_j\}_{j=1}^w)$ are obtained by minimizing a ridge-penalized least-squares problem~\cite{wasserman2006all}:
\begin{equation}\label{eq:mvar_ls_problem}
  \min_{A_0,\{A_j\}_{j=1}^w}\;
  \sum_{c=1}^{C}\;\sum_{k=w+1}^{K}
  \bigl\|y^{(c)}(t_k)-\hat y^{(c)}(t_k)\bigr\|_2^2.
\end{equation}

\noindent\textbf{Vectorized regression form.}
Define the stacked regressor
\[
  \phi^{(c)}_k \;:=\;
  \begin{bmatrix}1\\[2pt] y^{(c)}(t_{k-1})\\ \vdots\\ y^{(c)}(t_{k-w})\end{bmatrix}
  \!\in\mathbb{R}^{1+wr},
\]
and let
$B\in\mathbb{R}^{(1+wr)\times r}$ collect the intercept and lag matrices
so that $\hat y^{(c)}(t_k) = (\phi^{(c)}_k)^\top B$.
Stacking all $N_{\mathrm{train}}=C(K-w)$ training pairs into matrices
$X\in\mathbb{R}^{N_{\mathrm{train}}\times(1+wr)}$ and
$Y\in\mathbb{R}^{N_{\mathrm{train}}\times r}$,
ridge-regularized least squares solves
\begin{equation}\label{eq:mvar_ridge}
  \hat B \;=\; \arg\min_{B}\;\|Y - XB\|_F^2
             \;+\;\alpha\,\|B_{2:}\|_F^2,
\end{equation}
where $B_{2:}$ excludes the intercept row from regularization and
$\alpha>0$.
The closed-form solution is
\begin{equation}\label{eq:ridge_closed_form}
  \hat B \;=\; (X^\top X + \alpha\,\Gamma)^{-1}\,X^\top Y,
  \qquad
  \Gamma = \mathrm{diag}(0,1,\dots,1).
\end{equation}
The intercept row is excluded from the penalty term,
matching~\eqref{eq:ridge_closed_form}.

\noindent\textbf{Baseline values.}
In the baseline systematic comparison, the MVAR uses lag
$w_{\mathrm{MVAR}}=5$ and ridge strength
$\alpha=10^{-4}$.  At
$\Delta t_{\mathrm{rom}}=0.12$\,s this lag window spans $0.6$\,s of
physical time.  This is a compromise between capturing short-term latent
dependence and keeping the regression well-conditioned; ridge
regularization provides additional robustness.  Ridge regularisation
improves numerical stability of the fit, but it does not by itself guarantee
dynamical stability of the closed-loop rollout, which is determined by the
companion spectrum (Section~\ref{subsec:mvar_rollout}).  When the LSTM lag is
varied in dedicated ablations, the evaluation protocol uses a common
forecast start determined by the largest enabled lag in the config.
A formal unrestricted-VAR AIC/BIC study and an empirical lag sweep are
reported in Section~\ref{sec:lag_selection}; the final lag is ultimately
chosen on the basis of that end-to-end empirical comparison.

\subsection{Closed-loop rollout}\label{subsec:mvar_rollout}
For long-horizon forecasting the MVAR model is used in \emph{closed-loop}
(autoregressive) mode: each one-step prediction is fed back as input for
subsequent steps.

\noindent\textbf{Rollout map.}
Let the state of the autoregression be
\[
  z_k :=
  \begin{bmatrix}
    y(t_k)\\ y(t_{k-1})\\ \vdots\\ y(t_{k-w+1})
  \end{bmatrix}\in\mathbb{R}^{wr}.
\]
The MVAR($w$) model induces an affine update
\begin{equation}\label{eq:companion_update}
  z_{k+1} = \mathcal{A}_{\mathrm{comp}}\, z_k + b_{\mathrm{comp}},
\end{equation}
where $\mathcal{A}_{\mathrm{comp}}\in\mathbb{R}^{wr\times wr}$ is the
companion matrix with explicit block form
\begin{equation}
  \mathcal{A}_{\mathrm{comp}}
  =
  \begin{bmatrix}
    A_1 & A_2 & \cdots & A_{w-1} & A_w \\
    I_r & 0   & \cdots & 0       & 0   \\
    0   & I_r & \cdots & 0       & 0   \\
    \vdots & & \ddots &         & \vdots \\
    0   & 0   & \cdots & I_r     & 0
  \end{bmatrix},
  \qquad
  b_{\mathrm{comp}}
  =
  \begin{bmatrix} A_0 \\ 0 \\ \vdots \\ 0 \end{bmatrix},
  \label{eq:companion_matrix}
\end{equation}
where $I_r\in\mathbb{R}^{r\times r}$ is the identity and
$A_0\in\mathbb{R}^r$ the intercept; the top block row contains
$(A_1,\dots,A_w)$ and the subdiagonal contains identity ``shift'' blocks.
Errors compound through recursion, so closed-loop metrics are the primary
evaluation tool (Chapter~\ref{ch:exp_protocol_metrics}).

\noindent\textbf{Restrict--evolve--lift forecast.}
In the full pipeline the autoregressive forecast operates on physical-space
densities:
\begin{equation}\label{eq:closed_loop_forecast}
  \hat x^{(c)}(t_k) \;=\;
  \mathcal{L}\!\Bigl(
    \Phi_{\mathrm{ROM}}\!\bigl(
      \mathcal{R}(\hat x^{(c)}(t_{k-1})),\dots,
      \mathcal{R}(\hat x^{(c)}(t_{k-w}))
    \bigr)
  \Bigr),
\end{equation}
where $\mathcal{R}$ and $\mathcal{L}$ are the restriction and lifting
operators of Chapter~\ref{ch:global_pod_pipeline}, and $\Phi_{\mathrm{ROM}}$ is
either the MVAR or LSTM time-stepper.
The forecast is initialized with the first $w$ true (restricted) snapshots
and iterated for all remaining time steps.

\noindent\textbf{Stability.}
Closed-loop performance depends on the spectral radius
$\rho(\mathcal{A}_{\mathrm{comp}}) < 1$,
which is a sufficient condition for bounded rollout of the homogeneous part.
Ridge regularization shrinks coefficients and therefore tends to improve
stability margins.

\noindent\textbf{Optional spectral-radius scaling.}
An optional post-fit stabilization step
uniformly scales down coefficient matrices when
$\rho(\mathcal{A}_{\mathrm{comp}})$ exceeds the threshold.
This should be regarded as a modified model rather than a harmless
stabilisation, because it perturbs the learned dynamics.
In preliminary experiments it degraded forecast quality by over-damping
transient dynamics, so it is \emph{disabled} in all experimental runs.

\noindent\textbf{Parameter count.}\label{subsec:mvar_identifiability}
Each of the $r$ output dimensions is regressed on $w$ lagged $r$-vectors
plus a bias, giving $p_{\mathrm{MVAR}} = r + w\,r^2$ free parameters.
With $r{=}19$ and $w{=}5$ this yields $p_{\mathrm{MVAR}} = 1{,}824$,
fitted from $N_{\mathrm{train}} = C(K{-}w)$ samples across all training
trajectories; ridge regularisation keeps the effective parameter-to-data
ratio manageable.

%% ================================================================
\section{LSTM latent surrogate}
\label{sec:lstm_surrogate}
%% ================================================================

The LSTM serves as the \emph{nonlinear} latent time-stepper in the EF-ROM
benchmark.
It is introduced here as a flexible nonlinear comparator on the same latent
representation, not as an end in itself.
Because the LSTM has roughly $116\times$ more parameters than MVAR
(Section~\ref{subsec:mvar_identifiability}), any LSTM advantage is hard to
attribute purely to nonlinear structure versus overfitting capacity; the
LSTM's role is therefore \emph{diagnostic} rather than competitive.
Specifically:
\begin{itemize}[nosep]
  \item MVAR failure suggests either instability or low-dimensional mismatch.
  \item LSTM improvement over MVAR would suggest genuine nonlinear latent memory.
  \item LSTM non-improvement despite massive extra capacity is positive evidence
        that nonlinear expressivity is not the main bottleneck.
\end{itemize}

\noindent\textbf{Sequence-to-one formulation.}
Let $\mathbf{z}_k\in\mathbb{R}^d$ denote the POD coefficient vector at
time $t_k$ ($d = r = 19$).
A $p$-step autoregressive surrogate learns the one-step map
\begin{equation}\label{eq:lstm_seq_problem}
  \mathbf{z}_{k} \;\approx\;
  f_\theta\!\bigl(\mathbf{z}_{k-1},\ldots,\mathbf{z}_{k-p}\bigr),
  \qquad k = p{+}1,\ldots,K,
\end{equation}
where $f_\theta$ is a parameterized nonlinear model.
LSTMs are preferred over vanilla RNNs because the gated cell state
mitigates vanishing/exploding gradients, which matters when the latent
series exhibits long transients or regime shifts.

\noindent\textbf{LSTM architecture.}
Given an input window
$\mathbf{x}_k = (\mathbf{z}_{k-p},\ldots,\mathbf{z}_{k-1})
\in\mathbb{R}^{p\times d}$,
the LSTM processes the sequence through forget, input, and output gates
that regulate information flow into the cell state~$\mathbf{c}_\tau$ and
hidden state~$\mathbf{h}_\tau$ (gate equations in
Eq.~\eqref{eq:lstm_gates}, Appendix~\ref{app:supplementary_figures}).
The final hidden state is passed through LayerNorm~\cite{ba2016layernorm} before a linear readout.

\noindent\textbf{Residual (delta) prediction.}
In all experiments, the network uses residual prediction, so
the output layer predicts a \emph{correction} $\Delta_k$ and the final
prediction is
\begin{equation}\label{eq:lstm_residual}
  \widehat{\mathbf{z}}_{k} \;=\;
  \mathbf{z}_{k-1} \;+\; \Delta_k,
  \qquad
  \Delta_k := \mathbf{W}_y\,\mathrm{LayerNorm}(\mathbf{h}_{p}) + \mathbf{b}_y.
\end{equation}
This residual connection anchors the prediction to the most recent input
state, which improves stability during closed-loop rollout.

\noindent\textbf{Training objective and configuration.}
The LSTM is trained via one-step MSE on sliding windows of
ground-truth latent states, optimised with Adam~\cite{kingma2014adam}; early stopping,
gradient clipping, and a
learning-rate scheduler are used during training.
Full hyperparameters are listed in
Table~\ref{tab:lstm_config} (Appendix~\ref{app:supplementary_figures}).

\noindent\textbf{Parameter count.}
Each LSTM layer of hidden dimension~$h$ contributes
$4(d_{\mathrm{in}}h + h^2) + 8h$ weights (four gates $\times$
input-to-hidden, hidden-to-hidden, and two biases).
With $d{=}r{=}19$, $h{=}128$, $L{=}2$ layers, plus LayerNorm ($2h$) and
a linear readout ($hd+d$), the total is
$p_{\mathrm{LSTM}} = 76{,}288 + 132{,}096 + 256 + 2{,}451 = 211{,}091$---roughly
$116\times$ the MVAR count of $1{,}824$, a ratio that will become relevant
when interpreting forecast quality in Chapter~\ref{ch:rom_results}.

\noindent\textbf{Closed-loop forecasting.}
After training, we forecast by recursively feeding predictions back into
the input window:
$(\mathbf{z}_{k-p},\ldots,\mathbf{z}_{k-1})
\mapsto \widehat{\mathbf{z}}_{k}
\mapsto (\mathbf{z}_{k-p+1},\ldots,\mathbf{z}_{k-1},\widehat{\mathbf{z}}_{k})
\mapsto\cdots$
This closed-loop evaluation matches the restrict--evolve--lift protocol of
\eqref{eq:closed_loop_forecast}.
Unlike MVAR, the LSTM has no eigenvalue-based stability certificate;
stability is monitored empirically via rollout diagnostics
(boundedness of POD coefficients and mass drift after lifting).

\noindent\textbf{Algorithmic summary.}
Algorithm~\ref{alg:lstm_rom}
(Appendix~\ref{app:detailed_algorithms}) summarises the LSTM ROM pipeline.

\noindent\textbf{Comparison to MVAR.}
MVAR offers fast closed-form training and transparent stability analysis
(companion-matrix spectral radius), while LSTM provides nonlinear
capacity at the cost of more hyperparameters and weaker a~priori rollout
guarantees.
In crowd-dynamics EF-ROM benchmarks, linear MVARs can match or outperform
LSTMs in predictive accuracy while being lower-complexity and more
interpretable~\cite{alvarez2025eqfree}.

The comparison in Chapter~\ref{ch:rom_results} should therefore be
interpreted primarily as evidence about the effective geometry of the
aligned latent dynamics---whether they are primarily linear but
potentially unstable, substantially nonlinear, or insufficiently
low-dimensional for either surrogate to capture over long
horizons---rather than as a leaderboard-style predictor comparison.
